\documentclass[../DoAn.tex]{subfiles}
\begin{document}
\begin{longtable}{l p{11cm}}
	\hline
   \textbf{Thuật ngữ}  & \textbf{Giải thích} \\ \hline 
	\textbf{Godot Engine} & Game engine mã nguồn mở được sử dụng để xây dựng trò chơi trong đồ án. \\
	\textbf{GDScript} & Ngôn ngữ kịch bản tích hợp trong Godot, dùng để viết logic ứng dụng. \\
	\textbf{Scene} & Đơn vị tổ chức giao diện hoặc đối tượng trong Godot, gồm một cây các node. \\
	\textbf{Node} & Thành phần cơ bản trong Godot, có thể là phần tử giao diện, đối tượng hiển thị hoặc node xử lý logic. \\
	\textbf{Signal} & Cơ chế sự kiện của Godot, cho phép một node phát thông báo để node khác xử lý. \\
	\textbf{Ô chữ} & Dạng trò chơi trong đó các từ được đặt trên lưới và có thể giao nhau tại các ký tự chung. \\
	\textbf{Heuristic} & Cách tiếp cận tìm lời giải gần đúng dựa trên tiêu chí kinh nghiệm, không đảm bảo tối ưu toàn cục nhưng thường hiệu quả trong thực tế. \\
	\textbf{Chuẩn hóa từ} & Quá trình đưa từ về dạng thống nhất để xử lý, ví dụ chuyển về chữ thường hoặc bỏ khoảng trắng khi xếp bảng. \\

    \hline
\end{longtable}

\end{document}
